\chapter{Alta Disponibilidad (HA), NoSQL y Arquitecturas Cloud Modernas}
\label{cap:alta_disponibilidad_nosql_cloud}

\section{Marco Conceptual y Te\'orico (Curr\'iculo Universitario)}

La arquitectura de datos contempor\'anea exige tolerancia a fallos, recuperaci\'on ante desastres y escalabilidad horizontal tanto para cargas estructuradas como semiestructuradas.

\subsection{El Teorema CAP y Bases NoSQL}
Formulado por Eric Brewer, establece que en un sistema distribuido con particionamiento de red (\textbf{Partition Tolerance - P}), es imposible garantizar simult\'aneamente:
\begin{itemize}
    \item \textbf{Consistencia Fuerte (C):} Todos los nodos devuelven la misma versi\'on m\'as reciente del dato en todo momento.
    \item \textbf{Disponibilidad (A):} Toda petici\'on no fallida recibe una respuesta garantizada.
\end{itemize}

\subsection{Arquitecturas de Alta Disponibilidad en RDBMS}
\begin{itemize}
    \item \textbf{Streaming Replication (PostgreSQL):} Transmisi\'on continua de registros WAL desde un nodo primario (\textit{Read/Write}) hacia r\'eplicas en espera (\textit{Standby / Read Only}) de forma s\'incrona o as\'incrona.
    \item \textbf{Grupos de Disponibilidad AlwaysOn (SQL Server):} Replicaci\'on a nivel de base de datos con conmutaci\'on autom\'atica por error (\textit{Automatic Failover}).
    \item \textbf{Connection Pooling (PgBouncer):} Gesti\'on y reutilizaci\'on de conexiones para evitar la saturaci\'on de procesos en el servidor.
\end{itemize}

---

\section{Casos de Estudio: Almacenamiento Semiestructurado con JSONB y GIN}

\begin{lstlisting}[style=sqlstyle, caption={Modelado h\'ibrido relacional/documental en PostgreSQL}]
-- Tabla h\'ibrida con metadatos no estructurados en formato binario JSONB
CREATE TABLE device_telemetry (
    telemetry_id BIGSERIAL PRIMARY KEY,
    device_id VARCHAR(50) NOT NULL,
    recorded_at TIMESTAMPTZ DEFAULT CURRENT_TIMESTAMP,
    metrics JSONB NOT NULL
);

-- \'indice GIN para b\'usquedas ultra-r\'apidas dentro del documento JSONB
CREATE INDEX idx_telemetry_metrics_gin ON device_telemetry USING GIN (metrics);

-- Consulta con operador de contenci\'on (@>)
SELECT
    device_id,
    recorded_at,
    metrics->>'temperature' AS temp_celsius,
    metrics->'network'->>'ip_address' AS client_ip
FROM device_telemetry
WHERE metrics @> '{"status": "CRITICAL", "sensor_type": "thermal"}';
\end{lstlisting}

---

\section{Taller de Consolidaci\'on del Cap\'itulo}

\begin{businessproblem}
Una plataforma de telemetr\'ia IoT vehicular requiere almacenar registros de diagn\'osticos con estructuras de datos variables seg\'un el fabricante, permitiendo indexar propiedades internas del JSONB y consultar m\'etricas de fallo en milisegundos.
\end{businessproblem}

\subsection{Soluci\'on T\'ecnica con \'indices de Expresi\'on y GIN}

\begin{lstlisting}[style=sqlstyle, caption={Indexaci\'on especializada sobre campos JSONB anidados}]
-- 1. Inserci\'on de telemetr\'ia vehicular
INSERT INTO device_telemetry (device_id, metrics)
VALUES (
    'VIN-98234-A',
    '{
        "status": "CRITICAL",
        "sensor_type": "thermal",
        "temperature": 105.4,
        "engine": {"rpm": 4200, "oil_pressure_psi": 18.2},
        "error_codes": ["P0300", "P0117"]
    }'::jsonb
);

-- 2. \'indice de expresi\'on B-Tree para valores num\'ericos extra\'idos
CREATE INDEX idx_telemetry_temp_num ON device_telemetry (
    ((metrics->>'temperature')::NUMERIC)
) WHERE metrics ? 'temperature';

-- 3. Consulta optimizada de alerta
SELECT
    device_id,
    recorded_at,
    (metrics->>'temperature')::NUMERIC AS temp
FROM device_telemetry
WHERE metrics ? 'temperature'
  AND (metrics->>'temperature')::NUMERIC > 100.0;
\end{lstlisting}

\begin{engtip}
\begin{itemize}
    \item \textbf{Ventaja de \texttt{JSONB} frente a \texttt{JSON}:} \texttt{JSONB} descompone y almacena el documento en formato binario pre-parseado, eliminando espacios redundantes y soportando indexaci\'on GIN que acelera consultas de contenci\'on en \'ordenes de magnitud comparado con b\'usquedas de texto.
\end{itemize}
\end{engtip}
